% Los campos se completan en la versión que revisen y firmen los interesados.
\newcommand{\bloqueFirmaTesis}[2]{%
  \begin{minipage}[t]{.45\linewidth}
    \raggedright\small
    \rule{\linewidth}{.4pt}\par
    \textbf{#1}\par
    #2\par
  \end{minipage}%
}
\newcommand{\firmasTesis}{%
  \par\noindent
  \bloqueFirmaTesis{\asesorTesis}{Asesor\\%
    ORCID: \href{https://orcid.org/\orcidAsesor}{\orcidAsesor}}%
  \hfill
  \bloqueFirmaTesis{\autorTesis}{Autor}%
  \par
}
\newcommand{\estadoDeclaracionTesis}{%
  \noindent\textit{Borrador para revisión y firma}\par
}

\preliminarTesis[\centering]{Declaración de originalidad}
\estadoDeclaracionTesis

Yo, \asesorTesis, en calidad de asesor de la tesis titulada
«\textbf{\tituloTesis}», desarrollada por \autorTesis,
consigno los datos del reporte de similitud y el resultado de su revisión.

\noindent Índice de similitud: \rule{1.3cm}{.4pt}\,\%\par
\noindent Fecha del reporte: \rule{4cm}{.4pt}\par
\noindent Resultado de la revisión de las coincidencias, las citas y las referencias:\par
\noindent\rule{\linewidth}{.4pt}\par
\noindent\rule{\linewidth}{.4pt}\par
\noindent Ciudad y fecha: \ciudadTesis, \rule{6cm}{.4pt}\par

\vfill
\firmasTesis
\vspace*{.19\textheight}

\preliminarTesis{Acreditación}
\estadoDeclaracionTesis

La tesis titulada «\textbf{\tituloTesis}» corresponde a
\autorTesis, quien la presenta para optar el grado de Maestro en
\programaTesis{} en la Escuela de Posgrado de la Universidad Privada
Antenor Orrego.

La firma del asesor en esta hoja acredita la revisión integral del
informe, la atención de las observaciones formuladas durante la asesoría y
la conformidad para presentar la tesis al Jurado Evaluador.

\noindent\ciudadTesis, \rule{6cm}{.4pt}\par

\vfill
\firmasTesis
\vspace*{.19\textheight}

\preliminarTesis[\centering]{Presentación}

\vspace{.08\textheight}

Señores miembros del Jurado Evaluador:

En el marco del Reglamento de Grados y Títulos de la Escuela de Posgrado
de la Universidad Privada Antenor Orrego, someto a su consideración la
tesis titulada «\textbf{\tituloTesis}» para optar el grado académico de
Maestro en \programaTesis.

Agradezco de antemano la atención que dediquen a la lectura y evaluación
del informe, así como las observaciones y aportes que permitan mejorar
el trabajo de investigación.

\noindent\ciudadTesis, \anioTesis.\par

\vfill
\noindent\autorTesis\par
\vspace*{.15\textheight}

\clearpage
\thispagestyle{plain}
\vspace*{\fill}
\noindent\hfill
\begin{minipage}{.88\linewidth}
  \raggedleft\itshape
  \textbf{Dedicatoria}\par
  \medskip

  A Dios, por darme fortaleza y perseverancia para continuar mi formación
  profesional.

  A mis hijos, por su cariño y por la paciencia con la que me acompañaron
  durante esta etapa. A ellos dedico el esfuerzo de este trabajo.
\end{minipage}
\par\vspace*{.16\textheight}

\clearpage
\thispagestyle{plain}
\noindent\hfill
\begin{minipage}{.88\linewidth}
  \raggedleft\itshape
  \textbf{Agradecimiento}\par
  \vspace{.08\textheight}

  Agradezco a Dios por la fortaleza para perseverar durante la maestría
  y a mis hijos por su cariño y comprensión a lo largo de este proceso.

  Agradezco a los docentes de la Maestría en \programaTesis{} por su
  enseñanza y a mi asesor, \asesorTesis, por su orientación en el
  desarrollo de la investigación. Mi gratitud también a quienes
  colaboraron durante esta etapa.
\end{minipage}
\par\vspace{.05\textheight}
\begin{center}
  \autorTesis
\end{center}


\preliminarTesis[\centering]{Resumen}
\preliminarIndice{RESUMEN}
\begin{center}
\textbf{\tituloTesis}
\end{center}
El estudio tuvo como objetivo determinar la relación entre la condición del drenaje vial y el desempeño funcional del pavimento en el corredor EMP.LI105 Pedregal--Censo--Cumbray, Simbal, 2026, mediante un enfoque cuantitativo de alcance correlacional y diseño no experimental transversal sobre 20 unidades estructura--sector, cuya caracterización comprendió la continuidad longitudinal libre (CL) y la abertura transversal libre (CT) como indicadores de drenaje, junto con el índice de condición del pavimento (PCI) y el índice internacional de rugosidad (IRI) como medidas funcionales. Para examinar las cuatro relaciones se calcularon correlaciones de Spearman, se ajustaron las tendencias espaciales de los rangos y se efectuaron 9\,999 aleatorizaciones espectrales de Moran, cuyos valores de significancia se sometieron a la corrección familiar de Holm. Las correlaciones convencionales fueron de 0,5940 para CL--PCI, $-$0,6105 para CL--IRI, 0,1714 para CT--PCI y $-$0,1053 para CT--IRI; después del ajuste espacial y de multiplicidad, únicamente CL--PCI conservó significancia ($\rho$ residual = 0,5773; $p$ de Holm = 0,0180), lo que delimitó la evidencia a una asociación específica entre continuidad longitudinal y condición superficial. Estos resultados corresponden a una base de escenarios formativos destinada a la práctica académica, cuya procedencia se distingue de una campaña de medición de campo acreditada.

\noindent\textbf{Palabras clave:} drenaje vial; continuidad longitudinal; abertura transversal; condición del pavimento; rugosidad longitudinal.

\preliminarTesis[\centering]{Abstract}
\preliminarIndice{ABSTRACT}
\begin{center}
\textbf{Relationship between road drainage condition and functional pavement performance on the EMP.LI105 Pedregal--Censo--Cumbray road segment, Simbal, 2026}
\end{center}
This study aimed to determine the relationship between road drainage condition and functional pavement performance along the EMP.LI105 Pedregal--Censo--Cumbray road segment in Simbal, 2026, using a quantitative, correlational, non-experimental, cross-sectional design covering 20 drainage structure--road sector units; drainage condition was represented by free longitudinal drainage continuity (CL) and free transverse opening (CT), whereas functional performance was characterized by the Pavement Condition Index (PCI) and the International Roughness Index (IRI). To examine the four relationships, Spearman rank correlations were calculated, linear spatial trends were removed from the ranks, and 9,999 Moran spectral randomizations were carried out, with the resulting significance values subjected to Holm adjustment. Conventional correlations were 0.5940 for CL--PCI, $-$0.6105 for CL--IRI, 0.1714 for CT--PCI, and $-$0.1053 for CT--IRI; after spatial and multiple-testing adjustments, only CL--PCI remained statistically significant (residual $\rho$ = 0.5773; Holm-adjusted $p$ = 0.0180), indicating that the supported association was specific to longitudinal drainage continuity and pavement surface condition. These results come from a formative scenario dataset developed for academic practice, a source distinct from an accredited field measurement campaign.

\noindent\textbf{Keywords:} road drainage; longitudinal drainage continuity; transverse opening; pavement condition; longitudinal roughness.
